% !TEX root = ../../main.tex
\section{Cơ sở dữ liệu quan hệ và mở rộng Vector}
\label{sec:postgresql_redis}

\subsection{Hệ quản trị cơ sở dữ liệu quan hệ đối tượng PostgreSQL}
\label{subsec:postgresql}
PostgreSQL là hệ quản trị cơ sở dữ liệu quan hệ đối tượng mã nguồn mở (Object-Relational Database Management System -- ORDBMS), nổi bật với độ bền vững dữ liệu, khả năng mở rộng linh hoạt và cơ chế giao dịch đạt chuẩn ACID (Atomicity, Consistency, Isolation, Durability) tuyệt đối.

Các đặc điểm kỹ thuật then chốt của PostgreSQL bao gồm:
\begin{itemize}
    \item Cơ chế kiểm soát tương tranh đa phiên bản (MVCC): Cho phép các thao tác đọc không chặn ghi và ghi không chặn đọc thông qua cơ chế cô lập ảnh chụp (Snapshot Isolation), duy trì thông lượng truy vấn cao ngay cả khi nhiều học giả cùng truy cập đồng thời.
    \item Hỗ trợ dữ liệu bán cấu trúc (JSONB): Lưu trữ và đánh chỉ mục nhị phân các cấu trúc dữ liệu JSON phức tạp với hiệu suất truy vấn tương đương các trường dữ liệu quan hệ thông thường.
    \item Kiến trúc mở rộng mạnh mẽ (Extensibility): Cho phép tích hợp các kiểu dữ liệu mới, hàm người dùng tự định nghĩa và các tiện ích mở rộng (extensions) cấp thấp chạy trực tiếp trong nhân cơ sở dữ liệu.
\end{itemize}

Trong hệ thống VeriScholar, PostgreSQL giữ vai trò là xương sống lưu trữ dữ liệu trung tâm, kết hợp cùng dịch vụ giao diện lập trình ứng dụng FastAPI và mô hình kiến trúc phân tầng cô lập miền nghiệp vụ bám sát quy chế và hướng dẫn thực hiện đồ án tốt nghiệp đại học, quản lý toàn bộ thực thể học thuật, lịch sử phiên nghiên cứu, siêu dữ liệu phân đoạn tài liệu và bản đồ băm định danh nội dung (CAS).

\subsection{Tiện ích mở rộng pgvector và cấu trúc chỉ mục HNSW}
\label{subsec:pgvector_hnsw}
\texttt{pgvector} là một tiện ích mở rộng mã nguồn mở dành riêng cho PostgreSQL, bổ sung kiểu dữ liệu \texttt{vector} và các phép toán đo lường khoảng cách không gian đa chiều (khoảng cách Euclid $L_2$, khoảng cách Cosine, và tích vô hướng Inner Product). Nhờ \texttt{pgvector}, PostgreSQL hoạt động như một hệ quản trị cơ sở dữ liệu tích hợp đa mô hình quan hệ -- Vector (Integrated Relational-Vector DBMS), cho phép lưu trữ và truy vấn trực tiếp các Vector Embedding ngữ nghĩa 1536 chiều ngay bên cạnh dữ liệu quan hệ chuẩn ACID.

Điểm đột phá trong việc tối ưu hóa tốc độ tìm kiếm láng giềng gần nhất xấp xỉ (Approximate Nearest Neighbor -- ANN) của \texttt{pgvector} là việc hiện thực hóa thuật toán chỉ mục HNSW (Hierarchical Navigable Small World):
\begin{itemize}
    \item Nguyên lý phân tầng đa lớp: HNSW tổ chức không gian Vector thành cấu trúc đồ thị phân tầng tương tự như danh sách liên kết bước nhảy (Skip List). Các tầng trên cùng chứa các liên kết xa với mật độ thưa giúp định vị nhanh vùng lân cận tổng quát; các tầng dưới có mật độ nút dày đặc dần cho phép thu hẹp chính xác không gian tìm kiếm.
    \item Độ phức tạp truy vấn kỳ vọng thực nghiệm $O(\log N)$: Thay vì phải tính toán khoảng cách tuần tự trên toàn bộ kho vector (Flat Search với độ phức tạp $O(N)$), thuật toán định tuyến tham lam (Greedy Routing) trên đồ thị thế giới nhỏ phân tầng của HNSW đạt độ phức tạp kỳ vọng thực nghiệm tỷ lệ logarit $O(\log N)$ theo số lượng đỉnh $N$ trên các tập dữ liệu thực tế.
    \item Cập nhật linh hoạt và kiểm soát độ thu hồi theo siêu tham số: Khác với thuật toán IVFFlat đòi hỏi phải huấn luyện lại tâm cụm khi phân phối dữ liệu biến động lớn, HNSW cho phép chèn mới Vector trực tuyến. Độ thu hồi (Recall) và tốc độ truy vấn của chỉ mục được điều chỉnh linh hoạt thông qua ba siêu tham số cốt lõi: số liên kết hai chiều tối đa trên mỗi nút (\texttt{m}), kích thước danh sách ứng viên động khi dựng đồ thị (\texttt{ef\_construction}) và kích thước danh sách ứng viên khi tìm kiếm (\texttt{ef\_search}).
\end{itemize}

Việc hợp nhất dữ liệu quan hệ và dữ liệu Vector Embedding trong cùng một hệ quản trị PostgreSQL giúp loại bỏ hoàn toàn sự cồng kềnh của kiến trúc đa cơ sở dữ liệu, đảm bảo tính nhất quán giao dịch ACID trong suốt quá trình người dùng lưu trữ hoặc xóa tài liệu.

\subsection{Hệ thống đệm dữ liệu bộ nhớ truy cập nhanh Redis}
\label{subsec:redis}
Redis là hệ thống lưu trữ cấu trúc dữ liệu trong bộ nhớ RAM (In-Memory Data Store) mã nguồn mở với hiệu năng cao, đóng vai trò bộ nhớ đệm tốc độ cao trong hệ thống.

Trong kiến trúc VeriScholar, Redis đảm nhiệm ba vai trò chính: bộ đệm danh sách đen Token (Token Blacklist) thu hồi phiên JWT khi đăng xuất, kiểm soát giới hạn tần suất yêu cầu (Rate Limiting) bảo vệ cổng suy luận AI, và lưu tạm kết quả định danh học thuật (DOI Cache) giảm tải cho mạng ngoài.

\subsection{Luận giải lựa chọn cơ sở dữ liệu Vector tích hợp}
\label{subsec:vector_db_justification}
Một câu hỏi mang tính then chốt trong thiết kế hệ thống là: Tại sao VeriScholar lựa chọn giải pháp PostgreSQL tích hợp tiện ích pgvector thay vì sử dụng các cơ sở dữ liệu Vector chuyên dụng như Milvus, Pinecone hay Qdrant? Bảng~\ref{tab:vector_db_comparison} phân tích đối chứng chi tiết giữa hai cách tiếp cận trên các tiêu chí kỹ thuật cốt lõi.

\begin{table}[!htbp]
\centering
\caption{So sánh đa tiêu chí giữa PostgreSQL pgvector và cơ sở dữ liệu Vector chuyên dụng}
\label{tab:vector_db_comparison}
\renewcommand{\arraystretch}{1.08}
\begin{tabularx}{\textwidth}{|>{\centering\arraybackslash\bfseries}m{3.0cm}|>{\raggedright\arraybackslash}X|>{\raggedright\arraybackslash}X|}
\hline
\textbf{Tiêu chí} & \centering\textbf{PostgreSQL + pgvector (Đề tài chọn)} & \centering\textbf{Cơ sở dữ liệu Vector chuyên dụng (Milvus, Pinecone, Qdrant)} \tabularnewline
\hline
\textbf{Giao dịch ACID} & Tuyệt đối: Cập nhật bản ghi quan hệ và Vector nhúng đồng thời trong một giao dịch duy nhất. & Nhất quán cuối cùng: Rủi ro dữ liệu mồ côi hoặc bất đồng bộ giữa hai hệ thống lưu trữ. \\
\hline
\textbf{Kiến trúc vận hành} & Tối giản: Vận hành một cụm duy nhất, không phụ thuộc cụm dịch vụ trung gian (etcd, MinIO). & Phức tạp: Đòi hỏi cụm phân tán riêng biệt kèm hạ tầng điều phối tốn kém tài nguyên. \\
\hline
\textbf{Truy vấn kết hợp} & Linh hoạt: Kết hợp trực tiếp bộ lọc quan hệ (\texttt{WHERE}) và tìm kiếm Vector trong cùng một câu lệnh SQL. & Phân mảnh: Đòi hỏi cơ chế lọc hai pha tách rời gây phát sinh độ trễ mạng và nén ứng viên. \\
\hline
\textbf{Sao lưu phục hồi} & Chuẩn hóa: Sao lưu toàn diện và đồng bộ qua cơ chế nhật ký WAL và công cụ \texttt{pg\_dump}. & Tách rời: Cần quy trình đồng bộ sao lưu riêng lẻ giữa cơ sở dữ liệu nghiệp vụ và kho Vector. \\
\hline
\textbf{Chi phí hạ tầng} & Tiết kiệm: Mã nguồn mở 100\%, tối ưu hóa trên cấu hình máy chủ đám mây Azure đơn lẻ. & Tốn kém: Chi phí bản quyền SaaS định kỳ cao hoặc yêu cầu phần cứng RAM/CPU vượt trội. \\
\hline
\end{tabularx}
\end{table}

Qua phân tích đối chứng, việc lựa chọn PostgreSQL kết hợp \texttt{pgvector} là quyết định kỹ thuật chuẩn mực, giúp VeriScholar vừa bảo đảm tính toàn vẹn dữ liệu học thuật tuyệt đối, vừa tối ưu hóa hiệu năng và chi phí hạ tầng trên môi trường máy chủ Azure.
